\chapter*{第四部分导读：把方法放回天文与物理案例}
\addcontentsline{toc}{chapter}{第四部分导读：把方法放回天文与物理案例}

Part IV 是全书从方法训练走向完整案例的第一轮装配。前面我们分别学习了数据处理和机器学习 workflow；从这里开始，章节不再只围绕某一个概念展开，而围绕一个接近真实科研任务的问题展开。

本部分选择的案例包括 Gaia HR 图与恒星分类、光度红移估计、光谱分类和星系形态分类。它们覆盖了表格、颜色特征、光谱特征和图像形态，也自然对应回归、分类、特征工程、评估诊断和科学可信度。

\section*{案例章和方法章的区别}

方法章通常先讲概念，再给一个小例子；案例章则相反，先面对一个完整问题，再决定哪些工具真的有用。一个好的案例章应同时回答：

\begin{itemize}
    \item 数据来自哪里，字段或特征有什么物理含义；
    \item 目标变量或标签是否可靠；
    \item baseline 是什么，为什么不能直接跳到复杂模型；
    \item 模型错误如何回到样本和图表中解释；
    \item 结论能支持什么科学判断，哪些地方还只是候选线索。
\end{itemize}

这也是 Part IV 的教学重点：让学生看到科研项目不是“算法套题”，而是一组相互制约的判断。

\section*{读案例时的共同检查}

每读一个案例，都建议按同一套问题检查：

\begin{enumerate}
    \item 这个案例的科学问题是什么；
    \item 原始观测量和派生特征分别是什么；
    \item 模型要预测的目标是什么，它是否等同于真实物理类别；
    \item baseline 结果说明了什么；
    \item 复杂模型带来了什么改进，又带来了什么解释负担；
    \item 最终报告中哪些图表是证据，哪些只是辅助展示。
\end{enumerate}

\begin{protip}[案例不是炫技，而是装配]
Part IV 的目标不是堆更多模型，而是训练装配能力：把数据理解、简单模型、评估诊断、图表解释和报告写作组合成一条可交付的小项目。
\end{protip}
